\documentclass[10pt, a4paper]{article}

%=============================================================================
% CV LATEX TEMPLATE
%=============================================================================
%
% HOW TO USE THIS TEMPLATE
% ========================
% 1. Fill in the CONFIGURATION section below with personal details.
% 2. Replace the example entries in the DOCUMENT BODY with real content.
% 3. Add or remove sections as needed (Publications, Projects, etc.).
% 4. Each section provides one commented-out example showing the syntax.
% 5. Compile with pdflatex (or lualatex/xelatex).
%
% LLM INSTRUCTIONS
% ================
% When generating a CV from this template:
% - Fill in ALL \newcommand values in the CONFIGURATION block.
% - Use the provided entry commands: \educationentry, \experienceentry,
%   \projectentry, \techentry. Their signatures are documented below.
% - Separate multiple entries within a section with \entrybreak.
% - Bold key metrics/technologies with \textbf{}.
% - Escape special LaTeX characters: % → \%, & → \&, # → \#, $ → \$.
% - For links use \href{URL}{display text}.
% - Optional sections (Publications, Website) can be uncommented as needed.
% - Keep bullet points concise, action-verb-led, and quantified where possible.
%
% ENTRY COMMAND SIGNATURES
% ========================
% \educationentry{dates}{institution}{degree}{bullet items}
% \experienceentry{location}{dates}{title}{company}{bullet items}
% \projectentry{link-or-label}{title}{bullet items}
% \techentry{category}{comma-separated items}
%
%=============================================================================

%=============================================================================
% CONFIGURATION — Edit these values to customize your CV
%=============================================================================
% Personal information
\newcommand{\myname}{Your Full Name}
\newcommand{\mylocation}{City, Country}
\newcommand{\myemail}{your.email@example.com}
\newcommand{\myphone}{+1 234 567 890}
\newcommand{\mylinkedin}{https://www.linkedin.com/in/your-profile/}
\newcommand{\mylinkedinname}{Your Name}
\newcommand{\mygithub}{https://github.com/your-username}
\newcommand{\mygithubname}{your-username}

% Optional: personal website (uncomment the website block in HEADER to use)
% \newcommand{\mywebsite}{https://yourwebsite.com}
% \newcommand{\mywebsitename}{yourwebsite.com}

% Metadata
\newcommand{\lastupdated}{January 2026}

%=============================================================================
% PACKAGES
%=============================================================================
% Page layout
\usepackage[
    ignoreheadfoot,
    top=2cm, bottom=2cm, left=2cm, right=2cm,
    footskip=1.0cm,
    % showframe % uncomment for debugging
]{geometry}

% Typography and formatting
\usepackage{titlesec}
\usepackage{tabularx}
\usepackage{array}
\usepackage{enumitem}
\usepackage{amsmath}
\usepackage{needspace}

% Colors
\usepackage[dvipsnames]{xcolor}
\definecolor{primaryColor}{RGB}{0, 79, 144}

% Icons
\usepackage{fontawesome5}

% Links and PDF metadata
\usepackage[
    pdftitle={\myname\ — CV},
    pdfauthor={\myname},
    pdfcreator={LaTeX CV Template},
    colorlinks=true,
    urlcolor=primaryColor
]{hyperref}
\usepackage{bookmark}

% Layout helpers
\usepackage[pscoord]{eso-pic}
\usepackage{calc}
\usepackage{lastpage}
\usepackage{changepage}
\usepackage{paracol}
\usepackage{ifthen}
\usepackage{iftex}

%=============================================================================
% PDF/A COMPATIBILITY
%=============================================================================
\ifPDFTeX
    \input{glyphtounicode}
    \pdfgentounicode=1
    \usepackage[utf8]{inputenc}
    \usepackage{lmodern}
\fi

%=============================================================================
% SPACING CONSTANTS — Adjust these to change spacing throughout
%=============================================================================
\newlength{\entryspacing}
\setlength{\entryspacing}{0.2cm}

\newlength{\highlightspacing}
\setlength{\highlightspacing}{0.10cm}

\newlength{\sectiontopspacing}
\setlength{\sectiontopspacing}{0.3cm}

\newlength{\sectionbottomspacing}
\setlength{\sectionbottomspacing}{0.2cm}

%=============================================================================
% PAGE SETTINGS
%=============================================================================
\AtBeginEnvironment{adjustwidth}{\partopsep0pt}
\pagestyle{empty}
\setcounter{secnumdepth}{0}
\setlength{\parindent}{0pt}
\setlength{\topskip}{0pt}
\setlength{\columnsep}{0cm}

% Custom footer
\makeatletter
\let\ps@customFooterStyle\ps@plain
\patchcmd{\ps@customFooterStyle}{\thepage}{
    \color{gray}\textit{\small \myname\ — Page \thepage{} of \pageref*{LastPage}}
}{}{} 
\makeatother
\pagestyle{customFooterStyle}

%=============================================================================
% SECTION FORMATTING
%=============================================================================
\titleformat{\section}{\needspace{4\baselineskip}\bfseries\large}{}{0pt}{}[\vspace{1pt}\titlerule]
\titlespacing{\section}{-1pt}{\sectiontopspacing}{\sectionbottomspacing}

\renewcommand\labelitemi{$\circ$}

%=============================================================================
% CUSTOM ENVIRONMENTS — Do not modify unless changing layout
%=============================================================================

% Highlights (bullet points)
\newenvironment{highlights}{
    \begin{itemize}[
        topsep=\highlightspacing,
        parsep=\highlightspacing,
        partopsep=0pt,
        itemsep=0pt,
        leftmargin=0.4cm + 10pt
    ]
}{
    \end{itemize}
}

\newenvironment{highlightsforbulletentries}{
    \begin{itemize}[
        topsep=\highlightspacing,
        parsep=\highlightspacing,
        partopsep=0pt,
        itemsep=0pt,
        leftmargin=10pt
    ]
}{
    \end{itemize}
}

% One column entry
\newenvironment{onecolentry}{
    \begin{adjustwidth}{0.2cm + 0.00001cm}{0.2cm + 0.00001cm}
}{
    \end{adjustwidth}
}

% Two column entry
\newenvironment{twocolentry}[2][]{
    \onecolentry
    \def\secondColumn{#2}
    \setcolumnwidth{\fill, 4.5cm}
    \begin{paracol}{2}
}{
    \switchcolumn \raggedleft \secondColumn
    \end{paracol}
    \endonecolentry
}

% Header environment
\newenvironment{header}{
    \setlength{\topsep}{0pt}\par\kern\topsep\centering\linespread{1.5}
}{
    \par\kern\topsep
}

%=============================================================================
% CUSTOM COMMANDS — Do not modify unless changing layout
%=============================================================================

% Entry spacing command
\newcommand{\entrybreak}{\vspace{\entryspacing}}

% Last updated text
\newcommand{\placelastupdatedtext}{
    \AddToShipoutPictureFG*{
        \put(
            \LenToUnit{\paperwidth-2cm-0.2cm+0.05cm},
            \LenToUnit{\paperheight-1.0cm}
        ){\vtop{{\null}\makebox[0pt][c]{
            \small\color{gray}\textit{Last updated in \lastupdated}\hspace{\widthof{Last updated in August 2025}}
        }}}
    }
}

% Save original href
\let\hrefWithoutArrow\href

% External link with arrow icon
\renewcommand{\href}[2]{\hrefWithoutArrow{#1}{\ifthenelse{\equal{#2}{}}{ }{#2 }\raisebox{.15ex}{\footnotesize \faExternalLink*}}}

%=============================================================================
% CV ENTRY COMMANDS — Use these to create consistent entries
%=============================================================================

% Education entry: \educationentry{dates}{institution}{degree}{highlights}
\newcommand{\educationentry}[4]{
    \begin{twocolentry}{\textit{#1}}
        \textbf{#2}
        
        \textit{#3}
    \end{twocolentry}
    \vspace{\highlightspacing}
    \begin{onecolentry}
        \begin{highlights}
            #4
        \end{highlights}
    \end{onecolentry}
}

% Experience entry: \experienceentry{location}{dates}{title}{company}{highlights}
\newcommand{\experienceentry}[5]{
    \begin{twocolentry}{
        \textit{#1}
        
        \textit{#2}
    }
        \textbf{#3}
        
        \textit{#4}
    \end{twocolentry}
    \vspace{\highlightspacing}
    \begin{onecolentry}
        \begin{highlights}
            #5
        \end{highlights}
    \end{onecolentry}
}

% Project entry: \projectentry{link-or-label}{title}{highlights}
\newcommand{\projectentry}[3]{
    \begin{twocolentry}{\textit{#1}}
        \textbf{#2}
    \end{twocolentry}
    \vspace{\highlightspacing}
    \begin{onecolentry}
        \begin{highlights}
            #3
        \end{highlights}
    \end{onecolentry}
}

% Technology entry: \techentry{category}{items}
\newcommand{\techentry}[2]{
    \begin{onecolentry}
        \textbf{#1:} #2
    \end{onecolentry}
}


%=============================================================================
%=============================================================================
%
%  ██████   ██████   ██████ ██    ██ ███    ███ ███████ ███    ██ ████████
%  ██   ██ ██    ██ ██      ██    ██ ████  ████ ██      ████   ██    ██
%  ██   ██ ██    ██ ██      ██    ██ ██ ████ ██ █████   ██ ██  ██    ██
%  ██   ██ ██    ██ ██      ██    ██ ██  ██  ██ ██      ██  ██ ██    ██
%  ██████   ██████   ██████  ██████  ██      ██ ███████ ██   ████    ██
%
%  Replace the example entries below with your own content.
%  Each section shows one example entry; duplicate and modify as needed.
%
%=============================================================================
%=============================================================================
\begin{document}

% Header separator setup (internal — do not modify)
\newcommand{\AND}{\unskip\cleaders\copy\ANDbox\hskip\wd\ANDbox\ignorespaces}
\newsavebox\ANDbox
\sbox\ANDbox{}

%-----------------------------------------------------------------------------
% HEADER — Contact info is pulled from CONFIGURATION above
%-----------------------------------------------------------------------------
\placelastupdatedtext
\begin{header}
    \textbf{\fontsize{24pt}{24pt}\selectfont \myname}
    
    \vspace{0.3cm}
    
    \normalsize
    % Location
    \mbox{{\color{black}\footnotesize\faMapMarker*}\hspace*{0.13cm}\mylocation}%
    \kern 0.25cm\AND\kern 0.25cm%
    % Email
    \mbox{\hrefWithoutArrow{mailto:\myemail}{\color{black}{\footnotesize\faEnvelope[regular]}\hspace*{0.13cm}\myemail}}%
    \kern 0.25cm\AND\kern 0.25cm%
    % Phone
    \mbox{\hrefWithoutArrow{tel:\myphone}{\color{black}{\footnotesize\faPhone*}\hspace*{0.13cm}\myphone}}%
    \kern 0.25cm\AND\kern 0.25cm%
    % Website (uncomment the two lines below and add \mywebsite / \mywebsitename commands to CONFIGURATION)
    %\mbox{\hrefWithoutArrow{\mywebsite}{\color{black}{\footnotesize\faLink}\hspace*{0.13cm}\mywebsitename}}%
    %\kern 0.25cm\AND\kern 0.25cm%
    % LinkedIn
    \mbox{\hrefWithoutArrow{\mylinkedin}{\color{black}{\footnotesize\faLinkedinIn}\hspace*{0.13cm}\mylinkedinname}}%
    \kern 0.25cm\AND\kern 0.25cm%
    % GitHub
    \mbox{\hrefWithoutArrow{\mygithub}{\color{black}{\footnotesize\faGithub}\hspace*{0.13cm}\mygithubname}}%
\end{header}

%-----------------------------------------------------------------------------
% EDUCATION
% Syntax: \educationentry{dates}{institution}{degree}{bullet items}
% Use \entrybreak between entries.
%-----------------------------------------------------------------------------
\section{Education}

% --- Example education entry (replace with your own) ---
\educationentry{Sept 2022 – Jun 2026}
    {University Name}
    {BSc Your Degree}
    {
        \item Notable achievement or grade
        \item Another highlight
    }

% \entrybreak
%
% \educationentry{Sept 2020 – Jun 2022}
%     {Another Institution}
%     {Previous Qualification}
%     {
%         \item Achievement or detail
%     }

%-----------------------------------------------------------------------------
% EXPERIENCE
% Syntax: \experienceentry{location}{dates}{title}{company}{bullet items}
% Use \entrybreak between entries.
%-----------------------------------------------------------------------------
\section{Experience}

% --- Example experience entry (replace with your own) ---
\experienceentry{City, Country}{Jun 2025 – Sep 2025}
    {Job Title}
    {Company Name}
    {
        \item Describe what you built or achieved with quantified impact
        \item Mention key technologies with \textbf{bold}: e.g.\ \textbf{Python}, \textbf{Docker}
        \item Include metrics where possible: improved X by \textbf{30\%}
    }

% \entrybreak
%
% \experienceentry{City, Country}{Jan 2024 – May 2024}
%     {Another Role}
%     {Another Company}
%     {
%         \item Bullet point describing responsibilities and impact
%     }

%-----------------------------------------------------------------------------
% PUBLICATIONS (optional — uncomment if needed)
% Uses \section*{} (unnumbered) to exclude from numbering.
%-----------------------------------------------------------------------------
% \section*{Publications}
%
% \begin{itemize}
%     \item Author Name, ``Paper Title,'' \textit{Journal/Conference}, Year.
% \end{itemize}

%-----------------------------------------------------------------------------
% TECHNOLOGIES
% Syntax: \techentry{category}{comma-separated items}
% Use \entrybreak between entries.
%-----------------------------------------------------------------------------
\section{Technologies}

\techentry{Languages}{Language1, Language2, Language3}

\entrybreak

\techentry{Frameworks \& Tools}{Tool1, Tool2, Tool3}

%-----------------------------------------------------------------------------
% PROJECTS
% Syntax: \projectentry{link-or-label}{title}{bullet items}
%   - link-or-label: use \href{URL}{display text} for a link,
%     or plain text like "Work in Progress", "Course project", etc.
% Use \entrybreak between entries.
%-----------------------------------------------------------------------------
\section{Projects}

% --- Example project entry (replace with your own) ---
\projectentry{\href{https://github.com/your-username/project}{Git repository}}
    {Project Name}
    {
        \item One-sentence description of what the project does and why
        \item Tools used: Technology1, Technology2, Technology3
    }

% \entrybreak
%
% \projectentry{Course project}
%     {Another Project}
%     {
%         \item Description of the project
%         \item Tools used: X, Y, Z
%     }

\end{document}
